\section{Verification as LLM Guardrails}
\begin{frame}{An LLM proposes a construction; verification checks it}
\begin{center}
\begin{tikzpicture}[node distance=5mm,box/.style={draw=Teal,rounded corners,align=center,minimum height=17mm,text width=26mm,font=\small}]
\node[box] (s) {Human-reviewed\\contract};
\node[box,right=of s] (l) {LLM proposes\\code + proof};
\node[box,right=of l] (v) {Verus checks\\obligations};
\node[box,right=of v] (r) {Review\\and tests};
\draw[-{Stealth},thick,Teal] (s)--(l);
\draw[-{Stealth},thick,Teal] (l)--(v);
\draw[-{Stealth},thick,Teal] (v)--(r);
\draw[-{Stealth},thick,Rust] (v.south) -- ++(0,-.6) -| node[pos=.25,below,font=\scriptsize]{diagnostics / revision} (l.south);
\end{tikzpicture}
\end{center}
\vspace{6mm}
\end{frame}

\begin{frame}{A concrete task: count equal pairs}
Given a list $s$, count pairs of positions carrying the same value:
\[
 E(s)=\bigl|\{(i,j):0\le i<j<|s|,\ s_i=s_j\}\bigr|.
\]
For $s=[4,4,4,9,9]$, there are $\binom{3}{2}+\binom{2}{2}=4$ pairs.
\begin{block}{The recursive specification in \texttt{equal\_pairs.rs}}
\[
 E([])=0,\qquad E(p\mathbin{+\!+}[v])=E(p)+|\{i:p_i=v\}|.
\]
Appending $v$ creates one pair with each earlier occurrence of $v$.
\end{block}
\end{frame}

\begin{frame}[fragile]{The Spec and Task}
\begin{lstlisting}[language={},basicstyle=\ttfamily\fontsize{6}{7}\selectfont]
pub open spec fn spec_equal_pairs(s: Seq<u64>) -> nat
    decreases s.len()
{
    if s.len() == 0 {
        0
    } else {
        let prefix = s.drop_last();

        spec_equal_pairs(prefix)
            + prefix.filter(|x: u64| x == s.last()).len()
    }
}

fn equal_pairs(s: &[u64]) -> (r: u64)
    requires
        spec_equal_pairs(s@) <= u64::MAX
    ensures
        r == spec_equal_pairs(s@)
{
    <LLM can write here>
}

\end{lstlisting}

\end{frame}

\begin{frame}{The guardrail has a defined reach}
\begin{columns}[T]
\column{.48\textwidth}
\begin{block}{Can establish, when specified}
\begin{itemize}
\item The result is an attained maximum.
\item Slice accesses stay in bounds.
\item The prefix invariant is preserved.
\end{itemize}
\end{block}
\column{.48\textwidth}
\begin{block}{Still needs separate evidence}
\begin{itemize}
\item The contract matches user intent.
\item External systems match the model.
\item The algorithm is efficient enough to use.
\end{itemize}
\end{block}
\end{columns}
\vspace{6mm}
\large Stronger acceptance criteria make generated code easier to evaluate.
\end{frame}

\begin{frame}{In Practice}
\begin{enumerate}
\item Write the contract before polishing the implementation.
\item Gradually add contract components for existing code.
\item Make verification a required check for LLM-generated changes.
\end{enumerate}
\vspace{6mm}
\end{frame}

\begin{frame}{References and demo files}
\small
\begin{itemize}
\item \href{https://github.com/matthewgrimes/verus-demo}{This deck along with examples}\par\scriptsize\url{https://github.com/matthewgrimes/verus-demo}\small
\item \href{https://verus-lang.github.io/verus/guide/}{Verus Tutorial and Reference}\par\scriptsize\url{https://verus-lang.github.io/verus/guide/}\small
\item \href{https://verus-lang.github.io/verus/guide/}{Verus Tutorial and Reference}\par\scriptsize\url{https://verus-lang.github.io/verus/guide/}\small
\item \href{https://github.com/verus-lang/verus}{Verus source and installation instructions}\par\scriptsize\url{https://github.com/verus-lang/verus}\small
\item \href{https://microsoft.github.io/z3guide/}{Microsoft's Z3 Guide}\par\scriptsize\url{https://microsoft.github.io/z3guide/}\small
\end{itemize}
\medskip
Local examples: \texttt{examples/max.rs}, \texttt{examples/max\_iter.rs},\par
\texttt{examples/max.smt2}, \texttt{examples/equal\_pairs.rs}.
\end{frame}
